% !TeX program = xelatex
%% =====================================================================
%%  ใบกิจกรรมในชั้นเรียน (Worksheet) บทที่ 1 : ตรรกศาสตร์ --- สัปดาห์ที่ 1
%%  สำหรับแจกนักศึกษาทำในห้อง (ไม่มีเฉลย มีช่องว่างให้เขียน)
%% =====================================================================
\documentclass[a4paper,12pt]{article}
\usepackage{fontspec}
\usepackage{amsmath,amssymb}
\usepackage[margin=2cm]{geometry}
\usepackage{enumitem}
\usepackage{array}
\usepackage{xcolor}

\XeTeXlinebreaklocale "th_TH"
\XeTeXlinebreakskip = 0pt plus 1pt
\setmainfont[Script=Thai,Scale=1.25]{TH Sarabun New}
\definecolor{navy}{HTML}{1F3A93}
\renewcommand{\arraystretch}{1.5}
\setlist{topsep=4pt,itemsep=6pt}

%% เส้นว่างให้เขียนคำตอบ (#1 = จำนวนบรรทัด)
\newcounter{bl}
\newcommand{\blank}[1]{\par\vspace{0.2em}%
  \setcounter{bl}{0}\loop\ifnum\value{bl}<#1\relax
  \noindent\textcolor{gray!60}{\hrulefill}\par\vspace{1.0em}\stepcounter{bl}\repeat}

\begin{document}
\thispagestyle{empty}

%% ---------- หัวกระดาษ ----------
\begin{center}
  {\Large\bfseries\color{navy} ใบกิจกรรมในชั้นเรียน บทที่ 1 : ตรรกศาสตร์ (สัปดาห์ที่ 1)}\\[0.2em]
  {\normalsize รายวิชา 4091606 คณิตศาสตร์สำหรับคอมพิวเตอร์}
\end{center}
\vspace{0.3em}
\noindent\begin{tabular}{@{}p{0.5\textwidth} p{0.22\textwidth} p{0.22\textwidth}@{}}
  ชื่อ-สกุล: \dotfill & รหัส: \dotfill & กลุ่ม: \dotfill
\end{tabular}
\hrule
\vspace{0.6em}
\noindent\textbf{คำชี้แจง:} ให้นักศึกษาทำกิจกรรมต่อไปนี้ลงในใบกิจกรรม \textbf{ภายในเวลาที่กำหนด}
แสดงวิธีคิดให้ชัดเจน แล้วร่วมเฉลยและอภิปรายพร้อมกันในชั้นเรียน
\vspace{0.6em}

%% ---------- กิจกรรม 1 ----------
\noindent\textbf{กิจกรรมที่ 1} \hfill \textit{(5 นาที)}\\
ข้อความใดเป็น \textbf{ประพจน์}? ถ้าเป็น ให้ระบุค่าความจริง (T/F) ถ้าไม่เป็น ให้บอกเหตุผล
\begin{enumerate}[label=\arabic*)]
  \item ดวงอาทิตย์ขึ้นทางทิศตะวันออก \hfill \dotfill
  \item $5 + 7 = 13$ \hfill \dotfill
  \item เธอหิวข้าวหรือยัง? \hfill \dotfill
  \item $x - 2 = 9$ \hfill \dotfill
  \item ห้ามสูบบุหรี่ในอาคาร \hfill \dotfill
  \item จำนวนเต็มทุกจำนวนหารด้วย 2 ลงตัว \hfill \dotfill
\end{enumerate}
\vspace{0.4em}

%% ---------- กิจกรรม 2 ----------
\noindent\textbf{กิจกรรมที่ 2} \hfill \textit{(7 นาที)}\\
จงเติมตารางค่าความจริงของ \quad $\neg p \land q$ \quad (เติมคอลัมน์ $\neg p$ ก่อน)
\begin{center}
\begin{tabular}{|c|c|c|c|}
  \hline
  \textbf{$p$} & \textbf{$q$} & \textbf{$\neg p$} & \textbf{$\neg p \land q$} \\ \hline
  T & T & \rule{2.2cm}{0pt} & \rule{2.2cm}{0pt} \\ \hline
  T & F &  &  \\ \hline
  F & T &  &  \\ \hline
  F & F &  &  \\ \hline
\end{tabular}
\end{center}
\vspace{0.4em}

%% ---------- กิจกรรม 3 ----------
\noindent\textbf{กิจกรรมที่ 3} \hfill \textit{(5 นาที)}\\
จงหาค่าความจริงของ \quad $(p \rightarrow q) \land (\neg q \lor r)$ \quad เมื่อ $p=\text{T},\ q=\text{F},\ r=\text{T}$
(แสดงวิธีแทนค่าทีละขั้น)
\blank{4}

%% ---------- กิจกรรม 4 ----------
\noindent\textbf{กิจกรรมที่ 4} \hfill \textit{(5 นาที)}\\
กำหนด $p=$ ``ฉันอ่านหนังสือ'', \ $q=$ ``ฉันสอบผ่าน'' จงเขียนเป็นสัญลักษณ์
\begin{enumerate}[label=\arabic*)]
  \item ถ้าฉันอ่านหนังสือ แล้วฉันจะสอบผ่าน \hfill \dotfill
  \item ฉันสอบผ่านก็ต่อเมื่อฉันอ่านหนังสือ \hfill \dotfill
  \item ฉันไม่อ่านหนังสือ และฉันสอบไม่ผ่าน \hfill \dotfill
  \item ฉันอ่านหนังสือ หรือ ฉันสอบไม่ผ่าน \hfill \dotfill
\end{enumerate}
\vspace{0.4em}

%% ---------- กิจกรรม 5 ----------
\noindent\textbf{กิจกรรมที่ 5 (อภิปรายกลุ่ม)} \hfill \textit{(8 นาที)}\\
ร้านกาแฟติดป้ายว่า ``ลูกค้าจะได้รับส่วนลด \textbf{ถ้า} เป็นสมาชิก \textbf{และ} ซื้อครบ 100 บาท''\\
ให้ $M=$ เป็นสมาชิก, \ $B=$ ซื้อครบ 100 บาท, \ $D=$ ได้รับส่วนลด
\begin{enumerate}[label=\arabic*)]
  \item เขียนเงื่อนไขการได้ส่วนลดเป็นสัญลักษณ์ \hfill \dotfill
  \item ถ้าลูกค้าเป็นสมาชิกแต่ซื้อ 80 บาท จะได้ส่วนลดหรือไม่ เพราะเหตุใด
  \blank{2}
  \item สร้างตารางค่าความจริงของเงื่อนไขนี้
\begin{center}
\begin{tabular}{|c|c|c|}
  \hline
  \textbf{$M$} & \textbf{$B$} & \textbf{$D \equiv M \land B$} \\ \hline
  T & T & \rule{3cm}{0pt} \\ \hline
  T & F &  \\ \hline
  F & T &  \\ \hline
  F & F &  \\ \hline
\end{tabular}
\end{center}
\end{enumerate}

\end{document}
